\section{Supports and annihilators}
\label{section-supp-and-ann}

\noindent
Some very basic definitions and lemmas.

\begin{definition}
\label{definition-support-module}
Let $R$ be a ring and let $M$ be an $R$-module.
The {\it support of $M$} is the set
$$
\text{Supp}(M)
=
\{
\mathfrak p \in \Spec(R)
\mid
M_{\mathfrak p} \not = 0
\}
$$
\end{definition}

\begin{lemma}
\label{lemma-support-zero}
\begin{slogan}
A module over a ring has empty support if and only if it is the trivial module.
\end{slogan}
Let $R$ be a ring. Let $M$ be an $R$-module. Then
$$
M = (0) \Leftrightarrow \text{Supp}(M) = \emptyset.
$$
\end{lemma}

\begin{proof}
Actually,
Lemma \ref{lemma-characterize-zero-local}
even shows that $\text{Supp}(M)$ always contains a maximal ideal
if $M$ is not zero.
\end{proof}

\begin{definition}
\label{definition-annihilator}
Let $R$ be a ring. Let $M$ be an $R$-module.
\begin{enumerate}
\item Given an element $m \in M$ the {\it annihilator of $m$}
is the ideal
$$
\text{Ann}_R(m) = \text{Ann}(m) = \{f \in R \mid fm = 0\}.
$$
\item The {\it annihilator of $M$}
is the ideal
$$
\text{Ann}_R(M) = \text{Ann}(M) = \{f \in R \mid fm = 0\ \forall m \in M\}.
$$
\end{enumerate}
\end{definition}

\begin{lemma}
\label{lemma-annihilator-flat-base-change}
Let $R \to S$ be a flat ring map. Let $M$ be an $R$-module and
$m \in M$. Then $\text{Ann}_R(m) S = \text{Ann}_S(m \otimes 1)$.
If $M$ is a finite $R$-module, then
$\text{Ann}_R(M) S = \text{Ann}_S(M \otimes_R S)$.
\end{lemma}

\begin{proof}
Set $I = \text{Ann}_R(m)$. By definition there is an exact sequence
$0 \to I \to R \to M$ where the map $R \to M$ sends $f$ to $fm$. Using
flatness we obtain an exact sequence
$0 \to I \otimes_R S \to S \to M \otimes_R S$ which proves the first
assertion. If $m_1, \ldots, m_n$ is a set of generators of $M$
then $\text{Ann}_R(M) = \bigcap \text{Ann}_R(m_i)$. Similarly
$\text{Ann}_S(M \otimes_R S) = \bigcap \text{Ann}_S(m_i \otimes 1)$.
Set $I_i = \text{Ann}_R(m_i)$. Then it suffices to show that
$\bigcap_{i = 1, \ldots, n} (I_i S) = (\bigcap_{i = 1, \ldots, n} I_i)S$.
This is Lemma \ref{lemma-flat-intersect-ideals}.
\end{proof}

\begin{lemma}
\label{lemma-support-closed}
Let $R$ be a ring and let $M$ be an $R$-module.
If $M$ is finite, then $\text{Supp}(M)$ is closed.
More precisely, if $I = \text{Ann}(M)$ is the annihilator of $M$, then
$V(I) = \text{Supp}(M)$.
\end{lemma}

\begin{proof}
We will show that $V(I) = \text{Supp}(M)$.

\medskip\noindent
Suppose $\mathfrak p \in \text{Supp}(M)$. Then $M_{\mathfrak p} \not = 0$.
Choose an element $m \in M$ whose image in $M_\mathfrak p$ is nonzero.
Then the annihilator of $m$ is contained in $\mathfrak p$ by construction
of the localization $M_\mathfrak p$. Hence
a fortiori $I = \text{Ann}(M)$ must be contained in $\mathfrak p$.

\medskip\noindent
Conversely, suppose that $\mathfrak p \not \in \text{Supp}(M)$.
Then $M_{\mathfrak p} = 0$.
Let $x_1, \ldots, x_r \in M$ be generators.
By Lemma \ref{lemma-localization-colimit} there exists
an $f \in R$, $f\not\in \mathfrak p$ such that
$x_i/1 = 0$ in $M_f$. Hence $f^{n_i} x_i = 0$ for some $n_i \geq 1$.
Hence $f^nM = 0$ for $n = \max\{n_i\}$ as desired.
\end{proof}

\begin{lemma}
\label{lemma-support-base-change}
Let $R \to R'$ be a ring map and let $M$ be a finite $R$-module.
Then $\text{Supp}(M \otimes_R R')$ is the inverse image of
$\text{Supp}(M)$.
\end{lemma}

\begin{proof}
Let $\mathfrak p \in \text{Supp}(M)$. By Nakayama's lemma
(Lemma \ref{lemma-NAK}) we see that
$$
M \otimes_R \kappa(\mathfrak p) = M_\mathfrak p/\mathfrak p M_\mathfrak p
$$
is a nonzero $\kappa(\mathfrak p)$ vector space.
Hence for every prime $\mathfrak p' \subset R'$ lying
over $\mathfrak p$ we see that
$$
(M \otimes_R R')_{\mathfrak p'}/\mathfrak p' (M \otimes_R R')_{\mathfrak p'} =
(M \otimes_R R') \otimes_{R'} \kappa(\mathfrak p') =
M \otimes_R \kappa(\mathfrak p) \otimes_{\kappa(\mathfrak p)}
\kappa(\mathfrak p')
$$
is nonzero. This implies $\mathfrak p' \in \text{Supp}(M \otimes_R R')$.
For the converse, if $\mathfrak p' \subset R'$ is a prime lying
over an arbitrary prime $\mathfrak p \subset R$, then
$$
(M \otimes_R R')_{\mathfrak p'} =
M_\mathfrak p \otimes_{R_\mathfrak p} R'_{\mathfrak p'}.
$$
Hence if $\mathfrak p' \in \text{Supp}(M \otimes_R R')$
lies over the prime $\mathfrak p \subset R$, then
$\mathfrak p \in \text{Supp}(M)$.
\end{proof}

\begin{lemma}
\label{lemma-support-element}
Let $R$ be a ring, let $M$ be an $R$-module, and let $m \in M$.
Then $\mathfrak p \in V(\text{Ann}(m))$ if and only if
$m$ does not map to zero in $M_\mathfrak p$.
\end{lemma}

\begin{proof}
We may replace $M$ by $Rm \subset M$. Then (1) $\text{Ann}(m) = \text{Ann}(M)$
and (2) $m$ does not map to zero in $M_\mathfrak p$ if and only if
$\mathfrak p \in \text{Supp}(M)$.
The result now follows from Lemma \ref{lemma-support-closed}.
\end{proof}

\begin{lemma}
\label{lemma-support-finite-presentation-constructible}
Let $R$ be a ring and let $M$ be an $R$-module.
If $M$ is a finitely presented $R$-module, then $\text{Supp}(M)$ is a
closed subset of $\Spec(R)$ whose complement is quasi-compact.
\end{lemma}

\begin{proof}
Choose a presentation
$$
R^{\oplus m} \longrightarrow R^{\oplus n} \longrightarrow M \to 0
$$
Let $A \in \text{Mat}(n \times m, R)$ be the matrix of the first
map. By Nakayama's Lemma \ref{lemma-NAK} we see that
$$
M_{\mathfrak p} \not = 0 \Leftrightarrow
M \otimes \kappa(\mathfrak p) \not = 0 \Leftrightarrow
\text{rank}(A \bmod \mathfrak p) < n.
$$
Hence, if $I$ is the ideal of $R$ generated by the $n \times n$ minors
of $A$, then $\text{Supp}(M) = V(I)$. Since $I$
is finitely generated, say $I = (f_1, \ldots, f_t)$,
we see that $\Spec(R) \setminus V(I)$ is
a finite union of the standard opens $D(f_i)$, hence quasi-compact.
\end{proof}

\begin{lemma}
\label{lemma-support-quotient}
Let $R$ be a ring and let $M$ be an $R$-module.
\begin{enumerate}
\item If $M$ is finite then the support
of $M/IM$ is $\text{Supp}(M) \cap V(I)$.
\item If $N \subset M$, then $\text{Supp}(N) \subset
\text{Supp}(M)$.
\item If $Q$ is a quotient module of $M$ then $\text{Supp}(Q) \subset
\text{Supp}(M)$.
\item If $0 \to N \to M \to Q \to 0$ is a short exact sequence
then $\text{Supp}(M) = \text{Supp}(Q) \cup \text{Supp}(N)$.
\end{enumerate}
\end{lemma}

\begin{proof}
The functors $M \mapsto M_{\mathfrak p}$ are exact. This immediately
implies all but the first assertion. For the first assertion
we need to show that $M_\mathfrak p \not = 0$ and
$I \subset \mathfrak p$ implies $(M/IM)_{\mathfrak p}
= M_\mathfrak p/IM_\mathfrak p \not = 0$. This follows
from Nakayama's Lemma \ref{lemma-NAK}.
\end{proof}





